\section{Results and discussion}
\label{sec:results}

\subsection{Corrected exciton and biexciton energies}
\label{sec:energy-results}

The correlation corrections and corrected energies are collected in
\cref{tab:energies} of the Appendix. The results in
\cref{sec:energy-results,sec:binding-results,sec:overlap-results} are
dimensionless: lengths are scaled
by $\rB$ and energies by the effective Rydberg of
\cref{eq:effective-units}. At fixed effective-mass ratio
$\eta=\meff/\mheff$, they are material-scaled and do not require a dimensional
GaAs calibration. Conversion to dimensional units is used only for the
radiative-lifetime calculation. The values include the single-particle
confinement and correlation corrections but exclude the band gap. Because the
zeroth-order energies are analytical, each correlation correction has the
same numerical integration error as its corresponding corrected energy.
Although these rough errors are retained in the table, they are too small to
be visible on the scale of the biexciton-energy curves in \cref{fig:exx}.
The numerical values are reported for transparency and reproducibility, but
the physically relevant conclusions are the comparative geometry trends
within the idealized model rather than absolute predictions for a particular
quantum dot.

\begin{figure}[ht]
\centering
\includegraphics[width=\linewidth]{Main-E_XX-Ry.pdf}
\caption{Corrected biexciton energy as a function of
(a) the axial semiaxis $c$ at fixed $a=5\rB$ and (b) the lateral semiaxis
$a$ at fixed $c=\rB$, in effective Rydbergs. The band gap is excluded. Lines
are guides to the eye, and numerical error bars are smaller than the symbols
at the displayed scales.}
\label{fig:exx}
\end{figure}

Along the axial scan, $\EXX$ decreases steeply and nonlinearly, with most of
the change occurring between the two most strongly confined geometries. This
behavior reflects the strong sensitivity of the confinement energy to the
short semiaxis in the oblate regime.

The lateral scan is much flatter: enlarging the in-plane semiaxis produces
only a modest reduction in $\EXX$. The axial dimension is therefore the
dominant geometrical control of the absolute biexciton energy over the sampled
range, consistent with the separation of axial and lateral energy scales in
\cref{eq:single-particle-scales}.

\subsection{Biexciton binding}
\label{sec:binding-results}

\begin{figure}[ht]
\centering
\includegraphics[width=\linewidth]{Main-Binding-Energy-Ry.pdf}
\caption{Binding energy $\Ebind=2\EX-\EXX$ as a function of (a) $c$ at
fixed $a=5\rB$ and (b) $a$ at fixed $c=\rB$, in effective Rydbergs. Positive
values correspond to a bound biexciton. The dashed line marks zero binding.
Error bars are rough first-order propagations of the adaptive-integration
estimates, not statistical confidence intervals.}
\label{fig:binding}
\end{figure}

The central binding energy is positive throughout the sampled geometries and
decreases monotonically when either semiaxis is increased
(\cref{fig:binding}). Thus relaxing axial or lateral confinement reduces the
net correlation advantage of the four-particle state relative to two isolated
excitons. The lateral scan approaches the zero-binding threshold more closely
than the axial scan.

The binding is well separated from zero at every geometry except the widest
lateral dot. There, the numerical-resolution estimate is comparable to the
small central binding energy. Because that estimate omits systematic
integration bias and uncertainty in the optimum, we classify the point as
marginally positive rather than claiming firmly resolved binding. No sign
change is demonstrated within the studied range.

Because the single-particle part cancels in \cref{eq:binding-energy}, the
origin of this trend is exposed by comparing the two correlation corrections
in \cref{fig:correlation-corrections}. Both become less negative as the dot is
enlarged, while the separation between $2\Delta E_X$ and $\Delta E_{XX}$
narrows. That separation is precisely $\Ebind$, so the convergence of the two
curves accounts directly for the suppression of binding, particularly along
the lateral scan.

\begin{figure}[ht]
\centering
\includegraphics[width=\linewidth]{Main-Correlation-Corrections-Ry.pdf}
\caption{Comparison of the two-exciton correction $2\Delta E_X$ with the
biexciton correction $\Delta E_{XX}$ as functions of (a) $c$ at fixed
$a=5\rB$ and (b) $a$ at fixed $c=\rB$, in effective Rydbergs. Their vertical
separation equals the binding energy. Lines are guides to the eye; numerical
error bars are smaller than the plotted symbols at this scale.}
\label{fig:correlation-corrections}
\end{figure}

\FloatBarrier
\subsection{Variational parameters and optical intermediates}
\label{sec:overlap-results}

The complete optimized biexciton parameters are reported in
\cref{tab:parameters} of the Appendix and summarized graphically in
\cref{fig:parameters}. The exciton parameter $\alpha_X$ generally decreases along the
axial scan and increases along the lateral scan. For the biexciton, the
optimized cross-pair exponent $\beta$ becomes very small at the larger axial
sizes. The composite quantity $\gamma/\delta$, which sets the maximum of the isolated factor
$r_{ab}^{\gamma}e^{-\delta r_{ab}}$, increases smoothly as either dot
dimension is enlarged. This indicates that the hole--hole factor shifts
toward larger characteristic separations in larger dots.

\begin{figure}[ht]
\centering
\includegraphics[width=\linewidth]{Main-Variational-Parameters.pdf}
\caption{Optimized variational parameters along the axial scan at fixed
$a=5\rB$ and the lateral scan at fixed $c=\rB$. Panels (a) and (b) show
$\alpha_X$, $\alpha$, and $\beta$; panels (c) and (d) show $\gamma$ and
$\delta$. Lines are guides to the eye.}
\label{fig:parameters}
\end{figure}

\FloatBarrier

These parameter-level trends should be interpreted more cautiously than the
energies. The objective is numerically rough, the second-phase search is local,
and different parameter combinations can produce energies separated by less
than the practical quadrature resolution. The parameters characterize the
retained variational minimum but are not claimed to be unique.

The coincidence overlap is shown in \cref{fig:overlap}, while the derived
optical quantities are listed in \cref{tab:optical} of the Appendix. The
overlap increases as either dot dimension is enlarged. It tends toward
saturation at the larger axial sizes, whereas it continues to rise across the
lateral scan. The oscillator strength follows the same behavior.

\begin{figure}[ht]
\centering
\includegraphics[width=\linewidth]{Main-Overlap.pdf}
\caption{Squared exciton coincidence overlap $|\MX|^2$ as a function of
(a) $c$ at fixed $a=5\rB$ and (b) $a$ at fixed $c=\rB$. Lines are guides to
the eye, and numerical error bars are smaller than the plotted symbols.}
\label{fig:overlap}
\end{figure}

\FloatBarrier

\FloatBarrier
\subsection{Radiative lifetimes}
\label{sec:lifetime-results}

\begin{figure}[ht]
\centering
\includegraphics[width=\linewidth]{Main-Lifetimes.pdf}
\caption{Exciton and biexciton radiative lifetimes as functions of
(a) $c$ at fixed $a=5\rB$ and (b) $a$ at fixed $c=\rB$. The values use the
corrected optical transition energy and the coincidence overlap. The adopted
decay-channel convention gives $\tauXX=\tauX/4$. The single GaAs parameter
set is specified in \cref{sec:observables}. Propagated error bars are smaller
than the plotted symbols.}
\label{fig:lifetimes}
\end{figure}

Both radiative lifetimes decrease as the dot is enlarged
(\cref{fig:lifetimes}). Along the axial scan, the decrease weakens and becomes
nearly saturated for the largest dots. Along the lateral scan, the decline
remains pronounced across the sampled range. The biexciton lifetime follows
the same geometry dependence at one quarter of the exciton value.

Relaxing confinement decreases $E_{\mathrm{exc}}$, which alone would tend to
increase the lifetime. The simultaneous increase in $|\MX|^2$ is stronger,
however. Since substitution of \cref{eq:oscillator-strength} into
\cref{eq:lifetimes} gives
$\tauX\propto(E_{\mathrm{exc}}|\MX|^2)^{-1}$, the overlap enhancement
dominates and the lifetime decreases.

\FloatBarrier
\subsection{Comparison and limitations}
\label{sec:limitations}

The present geometry trends are qualitatively consistent with calculations
for full strongly oblate ellipsoidal dots, where the semiaxes tune the
correlated energies and binding and the exciton overlap sets the radiative
scale \cite{hayrapetyan2019,bleyan2022}. A direct numerical comparison is not
appropriate, however, because the semi-ellipsoidal boundary changes the axial
single-particle state, and the present four-particle trial function and full
integration protocol differ from those earlier calculations.

The model assumes impenetrable barriers, bulk effective masses, static
dielectric screening for the energy calculation, and an adiabatic
single-particle ground state. Finite band offsets, dielectric mismatch,
strain, atomistic symmetry, and phonon-assisted decay are outside the present
scope. The lifetime calculation follows the optical model and decay-channel
convention of Ref.~\cite{hayrapetyan2019}, using the single GaAs parameter set
specified in \cref{sec:observables}. It should therefore be interpreted as the
radiative lifetime within that model rather than as a complete experimental
decay time.

These approximations can shift the absolute energies, binding energies, and
lifetimes of a real structure. The principal outcome is therefore the
qualitative separation of axial and lateral geometry trends within a single
consistent model, not a quantitative prediction for a specific fabricated
dot.

Numerically, the reported bars quantify only first-order propagation of the
adaptive integrator's internal errors. The smooth trends and the positive
binding away from the widest lateral dot are robust on that scale, but a
higher-resolution targeted calculation would be required to determine firmly
whether the marginal endpoint is bound or antibound.

\FloatBarrier
